\chapter{Orientation Estimation and Sensor Fusion}

In the last two chapters we met the three little sensors that live inside an inertial
measurement unit (IMU) -- the accelerometer that feels gravity and motion, the
gyroscope that feels rotation, the magnetometer that feels the Earth's field -- and we
learned to calibrate the raw numbers into trustworthy physical units. But raw numbers are
not what we ultimately want. When you tilt your wrist to bend a note, or roll your
forearm to sweep a filter, the thing that matters is not ``what is the acceleration
along the chip's $x$-axis right now'' -- it is \emph{which way is the wrist pointing}.
That single question -- the device's \emph{orientation} in space, also called its
\emph{attitude} or \emph{pose} -- is the subject of this chapter.

Orientation is a deceptively rich quantity. It is not three numbers you can simply add
up; it lives on a curved space (the set of rotations), and the naive ways of writing
it down hide nasty traps. Worse, none of our three sensors measures orientation
directly. The gyro measures the \emph{rate of change} of orientation, so to get
orientation you must integrate -- and integration quietly accumulates error until your
estimate slowly walks away from the truth. The accelerometer and magnetometer, by
contrast, give you \emph{absolute} references (down, and roughly north) but they are
noisy and, for the accelerometer, corrupted by every real motion. Neither sensor alone
is trustworthy. The art of combining them so that each covers the other's weakness is
called \emph{sensor fusion}, and it is one of the most satisfying applications of the
ideas in the companion linear-algebra and probability books: orientation as geometry,
fusion as optimal estimation under noise.

We will build the picture from the ground up. First, how do you even \emph{represent}
a rotation -- Euler angles and their gimbal-lock trap, rotation matrices, and the
quaternion that practitioners actually use? Then, why does the gyro drift and why are
the accelerometer and magnetometer ``absolute but noisy''? Then the central trick: the
\emph{complementary filter}, which is just a high-pass on the gyro plus a low-pass on
the accelerometer, fused in the frequency domain you already understand. We will derive
the \emph{Kalman filter} intuitively as the optimal version of that same idea, and meet
the \emph{Madgwick} and \emph{Mahony} filters that run on real microcontrollers. We
close, as always, with the wrist tracker: how to get a stable wrist orientation, and a
clear accounting of what orientation can and cannot tell you about the hand.

\section{What orientation is, geometrically}
\label{ch36:what-orientation-is}

Fix two coordinate frames. One is glued to the world: call it the \emph{navigation
frame} or \emph{Earth frame}, with axes pointing, say, North, East, and Down (the NED
convention). The other is glued to the sensor chip and tumbles with it: the \emph{body
frame}, with axes along the chip's own $x$, $y$, $z$. \emph{Orientation} is the answer
to a single question: how is the body frame rotated relative to the world frame? Every
way of writing down orientation is just a different bookkeeping for that one rotation.

\begin{intuitionbox}
Hold your phone flat on the table, screen up. Body frame and world frame agree --
zero rotation. Now tip the phone up to read it: you have rotated the body frame about
a horizontal axis. Spin in your chair while holding it: another rotation, about
vertical. Orientation is the \emph{total} rotation that carries the world axes onto the
body axes. The point of this chapter is that ``a rotation'' is a single geometric
object, and Euler angles, matrices, and quaternions are three languages for naming the
same object.
\end{intuitionbox}

A crucial fact, easy to forget, is that rotations do \emph{not commute}. Tip then
spin lands you somewhere different from spin then tip. Try it with the phone: the final
orientations genuinely differ. This non-commutativity is the deep reason rotations are
subtle, and it is exactly what gimbal lock and quaternion algebra are wrestling with.

\subsection{Euler angles: three turns in a row}
\label{ch36:euler}

The most human way to describe an orientation is as a sequence of three simple turns
about axes: \emph{roll} $\phi$ about the forward $x$-axis, \emph{pitch} $\theta$ about
the side $y$-axis, and \emph{yaw} $\psi$ about the down $z$-axis. Aircraft people have
used these names for a century: roll dips a wingtip, pitch raises the nose, yaw swings
the nose left or right. Any orientation can be reached by some $(\phi,\theta,\psi)$.

\begin{figure}[H]
\centering
\begin{tikzpicture}[scale=1.0,>=Stealth]
  % body axes
  \draw[->,thick,sourceblue] (0,0,0) -- (2.4,0,0) node[right]{$x$ (roll, $\phi$)};
  \draw[->,thick,targetgreen] (0,0,0) -- (0,2.0,0) node[above]{$z$ (yaw, $\psi$)};
  \draw[->,thick,quotientorange] (0,0,0) -- (0,0,2.4) node[below left]{$y$ (pitch, $\theta$)};
  % rotation arcs
  \draw[->,sourceblue,thick] (1.6,0,0) ++(0,0.5,0) arc (90:-30:0.5);
  \draw[->,quotientorange,thick] (0,0,1.6) ++(0.5,0,0) arc (0:120:0.5);
  \node[structurepurple] at (1.9,1.4,0) {\small body frame};
\end{tikzpicture}
\caption{Euler angles describe orientation as three successive turns: roll about the
forward axis, pitch about the side axis, yaw about the vertical axis. Intuitive, but
the order matters and the scheme has a singularity (Section~\ref{ch36:gimbal}).}
\label{ch36:fig-euler}
\end{figure}

Euler angles are wonderful for talking to humans -- ``pitch up thirty degrees'' is
instantly clear -- and they are how we will report the final answer. But three numbers
hide three problems. First, the \emph{order} matters (roll-pitch-yaw $\neq$
yaw-pitch-roll), so you must commit to a convention. Second, the angles \emph{wrap}:
$359^\circ$ and $-1^\circ$ are the same heading, so naive averaging or differencing
goes wrong near the wrap. Third, and most dangerous, is gimbal lock.

\subsection{The gimbal-lock trap}
\label{ch36:gimbal}

\begin{trapbox}
When pitch reaches $\pm 90^\circ$, the roll axis and the yaw axis line up and become
the \emph{same} physical axis. You have lost a degree of freedom: two of your three
numbers now do the same thing, and the third has become undefined. This is
\emph{gimbal lock}. Near it, tiny real motions demand huge, fast changes in the Euler
angles, so any estimator that stores orientation as Euler angles becomes numerically
violent there.
\end{trapbox}

\begin{intuitionbox}
Picture a camera on a mechanical gimbal with three nested rings: outer ring yaws,
middle ring pitches, inner ring rolls. Pitch the middle ring up to $90^\circ$ so the
camera stares at the sky. Now the inner (roll) ring and the outer (yaw) ring spin about
the \emph{same} vertical line. Spinning either one just rotates the picture; you can no
longer tip the camera sideways without moving the locked middle ring. A whole direction
of motion has vanished from your control. The Apollo guidance computer had a real
gimbal-locked inertial platform; astronauts were warned to avoid certain attitudes, and
mission control joked grimly about ``going to gimbal lock.''
\end{intuitionbox}

Gimbal lock is not a bug you can patch out of Euler angles; it is a topological fact.
There is no way to cover the full space of 3-D rotations with three smooth angle
coordinates without a singularity somewhere -- the same reason every flat map of the
globe distorts the poles. The fix is to stop forcing orientation into three numbers.
We move to representations that are \emph{global} and \emph{smooth}: rotation matrices
(nine numbers, six constraints) and quaternions (four numbers, one constraint).

\subsection{Rotation matrices}
\label{ch36:rotmat}

A rotation in three dimensions is a linear map that preserves lengths and angles and
keeps the handedness of space. In the language of the linear-algebra companion, it is an
orthogonal matrix with determinant $+1$ -- an element of the group $SO(3)$. Write
$\mat{R}$ for the matrix that takes a vector's coordinates in the body frame and gives
its coordinates in the world frame:
\[
  \vec{v}_{\text{world}} = \mat{R}\,\vec{v}_{\text{body}}, \qquad
  \mat{R}^\top \mat{R} = \mat{I}, \qquad \det \mat{R} = +1 .
\]
The columns of $\mat{R}$ are the body axes expressed in world coordinates -- so $\mat R$
literally \emph{is} a picture of where the chip's axes are pointing. The single
elementary rotations are the ones you would guess; for example a rotation by $\theta$
about the $y$-axis is
\[
  \mat{R}_y(\theta) =
  \begin{bmatrix} \cos\theta & 0 & \sin\theta \\ 0 & 1 & 0 \\ -\sin\theta & 0 & \cos\theta \end{bmatrix},
\]
and the full Euler orientation is a product like
$\mat{R} = \mat{R}_z(\psi)\,\mat{R}_y(\theta)\,\mat{R}_x(\phi)$. The non-commutativity
of rotations is now obvious: matrix multiplication does not commute.

\begin{whybox}
Rotation matrices never gimbal-lock and compose by plain multiplication, which makes
them excellent for \emph{computing}. Their drawback is redundancy: nine numbers for
three degrees of freedom, tied by six constraints ($\mat R^\top\mat R=\mat I$). When you
integrate a matrix forward in time with finite-precision arithmetic, those constraints
slowly rot -- the columns drift out of orthonormality -- so you must periodically
\emph{re-orthonormalize}. Quaternions carry only one constraint and are cheaper to keep
clean, which is why they win in practice.
\end{whybox}

\subsection{Quaternions: the practitioner's choice}
\label{ch36:quat}

The quaternion is the representation almost every real attitude filter uses internally.
It is a four-number object $q = (q_w,\,q_x,\,q_y,\,q_z)$, often written
$q = q_w + q_x\jj_1 + q_y\jj_2 + q_z\jj_3$ with three imaginary units. The geometry
behind the algebra is clean and is the only part you must \emph{feel}:

\begin{intuitionbox}
\textbf{Euler's rotation theorem.} Any orientation, however complicated, equals a single
rotation by some angle $\alpha$ about some fixed axis $\hat{\vec n}$ (a unit vector).
A unit quaternion packs exactly that axis and angle into four numbers:
\[
  q = \Big(\cos\tfrac{\alpha}{2},\ \ \hat n_x \sin\tfrac{\alpha}{2},\ \
            \hat n_y \sin\tfrac{\alpha}{2},\ \ \hat n_z \sin\tfrac{\alpha}{2}\Big).
\]
The scalar part remembers \emph{how much} you turned; the vector part remembers
\emph{about what axis}. The half-angle is the famous quirk -- a full $360^\circ$ turn
gives $-q$ and a $720^\circ$ turn returns to the start. For our purposes you only need:
unit quaternions name rotations, and there are no singularities.
\end{intuitionbox}

Three facts let you use quaternions without drowning in algebra. (1) A \emph{unit}
quaternion, $\norm{q} = 1$, represents a rotation; you renormalize after each update to
keep it on the unit sphere -- one constraint, cheap. (2) Rotations \emph{compose} by
the quaternion product $q_2 \otimes q_1$ (do $q_1$ then $q_2$), which is just a fixed
bilinear formula. (3) To rotate a vector $\vec v$, sandwich it:
$\vec v' = q \otimes \vec v \otimes q^{-1}$. The quaternion product is non-commutative,
exactly mirroring the non-commutativity of rotations.

\begin{formalbox}
The Hamilton product of $a=(a_w,\vec a)$ and $b=(b_w,\vec b)$ (scalar part, vector part) is
\[
  a\otimes b = \big(a_w b_w - \vec a\cdot\vec b,\ \ a_w \vec b + b_w \vec a + \vec a\times\vec b\big).
\]
For a unit quaternion the inverse is the conjugate $q^{-1}=q^{*}=(q_w,-\vec q)$. The
rotation matrix equivalent to $q$ is
\[
  \mat R(q)=\begin{bmatrix}
  1-2(q_y^2+q_z^2) & 2(q_xq_y-q_wq_z) & 2(q_xq_z+q_wq_y)\\
  2(q_xq_y+q_wq_z) & 1-2(q_x^2+q_z^2) & 2(q_yq_z-q_wq_x)\\
  2(q_xq_z-q_wq_y) & 2(q_yq_z+q_wq_x) & 1-2(q_x^2+q_y^2)
  \end{bmatrix}.
\]
You convert to Euler angles only at the very end, for display, and you keep the
quaternion as the working state.
\end{formalbox}

\begin{historybox}
Quaternions were carved by William Rowan Hamilton into the stone of Broom Bridge in
Dublin on 16 October 1843, in a flash of insight while walking with his wife: the
relation $\jj_1^2=\jj_2^2=\jj_3^2=\jj_1\jj_2\jj_3=-1$. For decades they were a
mathematical curiosity, eclipsed by vector calculus. Their revenge came with the space
age and computer graphics, where their freedom from gimbal lock and their cheap,
stable composition made them indispensable. Every spacecraft, drone, and VR headset
today tracks attitude with quaternions.
\end{historybox}

\begin{connectionbox}
This is the linear-algebra companion made physical. Rotation matrices are the orthogonal
matrices ($\mat R^\top\mat R=\mat I$) from the chapter on orthonormal bases; their
columns are an orthonormal basis -- the body axes. Quaternions are the unit-norm
elements of a 4-D algebra; ``rotate by sandwiching'' is a conjugation, the same pattern
as a change of basis $\mat P \mat A \mat P^{-1}$. Orientation lives on a curved manifold
($SO(3)$), which is precisely why no flat three-number chart (Euler angles) can cover it
without a singularity.
\end{connectionbox}

\section{Why each sensor is not enough alone}
\label{ch36:why-not-enough}

We now have a language for orientation. Where do the numbers come from? Each sensor
offers a partial, flawed view. Fusion is worth doing only because the flaws are
\emph{complementary}.

\subsection{The gyroscope: smooth but drifting}
\label{ch36:gyro-drift}

A gyroscope measures \emph{angular rate} $\boldsymbol\omega(t)$ -- how fast the body is
rotating, in rad/s, about each body axis. To get orientation you integrate. In
quaternion form the kinematics are a clean differential equation,
\[
  \dot q = \tfrac{1}{2}\, q \otimes (0,\ \boldsymbol\omega),
\]
which in a microcontroller becomes the discrete step $q_{k} = q_{k-1} \otimes
\Delta q(\boldsymbol\omega_k,\Delta t)$, renormalized. Integration is the gyro's gift
and its curse. Because it accumulates rate, the gyro is remarkably \emph{smooth} and
responsive over short times -- it tracks fast wrist flicks accurately. But it has no
idea where ``down'' is, and every small error in the rate accumulates.

\begin{happeningbox}
Two things make integrated gyro orientation walk away from the truth. The big one is
\emph{bias}: a real gyro reads a small nonzero rate even when perfectly still (recall
the calibration chapter). Integrating a constant bias $b$ gives an angle error that
grows \emph{linearly}, $b\cdot t$ -- a steady drift. The second is \emph{random-walk}:
broadband noise on the rate, once integrated, performs a random walk whose spread grows
like $\sqrt{t}$. Bias dominates after a few seconds. A cheap MEMS gyro might drift a
degree or more per second -- tolerable for a one-second gesture, ruinous over a minute.
\end{happeningbox}

\begin{figure}[H]
\centering
\begin{tikzpicture}
\begin{axis}[width=11.5cm,height=5.0cm,
  xlabel={time $t$ (s)}, ylabel={heading error (deg)},
  xmin=0,xmax=30,ymin=-5,ymax=35, domain=0:30, samples=120,
  legend pos=north west, legend cell align=left, grid=both,
  axis lines=left, every axis plot/.append style={thick}]
  \addplot[sourceblue,thick] {1.0*x};
  \addlegendentry{integrated gyro (drift $\approx 1^\circ$/s)}
  \addplot[specred,only marks,mark=*,mark size=0.6pt,samples=120,domain=0:30]
      {2.2*sin(deg(2.3*x))+1.6*cos(deg(1.1*x+1))};
  \addlegendentry{accel/mag estimate (noisy, no drift)}
  \addplot[targetgreen,thick,dashed] {0.0};
  \addlegendentry{truth}
\end{axis}
\end{tikzpicture}
\caption{The complementary flaws. The integrated gyro (blue) is smooth but ramps away
from truth as bias accumulates. The accelerometer/magnetometer estimate (red dots) is
jittery but stays centered on truth -- no drift. Fusion keeps the gyro's smoothness at
short times and the accel/mag's long-term anchoring.}
\label{ch36:fig-drift}
\end{figure}

\subsection{Accelerometer and magnetometer: noisy but absolute}
\label{ch36:absolute}

A stationary accelerometer measures the reaction to gravity: it reports an upward $1g$
specific force, which points along world ``up'' regardless of how long you have been
running. That gives you two of the three orientation angles directly -- \emph{roll} and
\emph{pitch} (the tilt of the body relative to gravity) -- as
\[
  \phi = \operatorname{atan2}(a_y,\,a_z), \qquad
  \theta = \operatorname{atan2}\!\big(-a_x,\ \sqrt{a_y^2+a_z^2}\big).
\]
These are \emph{absolute}: they do not drift, because gravity does not drift. The
magnetometer plays the same role for the remaining angle, \emph{yaw} (heading): the
Earth's field points roughly toward magnetic north, giving an absolute compass bearing
once you have tilt-compensated it using the accel-derived roll and pitch.

\begin{trapbox}
The accelerometer measures gravity \emph{plus all other acceleration}. The moment the
wrist actually moves -- and a gesture is nothing but motion -- the accel vector is
gravity contaminated by the motion you are trying to track. So the accel's ``down'' is
only trustworthy when the device is nearly still, or on average over time. Likewise the
magnetometer is wrecked by nearby steel, motors, speakers, and the very electronics on
your board (hard- and soft-iron distortion from the calibration chapter). Both
absolute references are \emph{noisy and intermittently wrong}; only the gyro is reliable
moment to moment.
\end{trapbox}

The pattern is now sharp, and it is a frequency-domain pattern, which means we already
own the tools to exploit it:
\begin{itemize}
  \item The \textbf{gyro} is trustworthy at \emph{high} frequencies (fast changes) and
        untrustworthy at \emph{low} frequencies / DC (slow drift, bias).
  \item The \textbf{accel/mag} are trustworthy at \emph{low} frequencies (the long-term
        average really is down/north) and untrustworthy at \emph{high} frequencies
        (motion and noise).
\end{itemize}
Two signals that are each good in opposite halves of the spectrum -- and we know exactly
how to glue those: a high-pass on one, a low-pass on the other, added together so the
filters sum to unity. That is the complementary filter.

\section{The complementary filter}
\label{ch36:complementary}

\begin{whybox}
Let $\theta_g$ be the orientation you get by integrating the gyro and $\theta_a$ the
orientation you get from the accelerometer (and magnetometer). Choose a high-pass
$H(s)$ and a low-pass $L(s)$ that are \emph{complementary}: $H(s)+L(s)=1$ for every
frequency. Then the fused estimate
\[
  \hat\theta = H(s)\,\theta_g + L(s)\,\theta_a
\]
keeps the gyro's fast detail (high-pass kills the gyro's low-frequency drift) and the
accel/mag's slow truth (low-pass kills their high-frequency motion noise), and because
$H+L=1$ a steady true orientation passes through with unit gain and no distortion. You
are not throwing away either sensor; you are using each only where it is reliable.
\end{whybox}

The discrete-time form is one line of code and is what runs on a real board. With a
single tuning constant $\alpha$ close to $1$,
\[
  \boxed{\ \hat\theta_k = \alpha\,\big(\hat\theta_{k-1} + \omega_k\,\Delta t\big)
           \;+\; (1-\alpha)\,\theta_{a,k}\ }
\]
Read it geometrically: \emph{predict} the new angle by stepping the previous estimate
forward with the gyro ($\hat\theta_{k-1}+\omega_k\Delta t$), then nudge that prediction
a little ($1-\alpha$ of the way) toward what the accelerometer says. The $\alpha$ term
is the high-pass on the gyro; the $(1-\alpha)$ term is the low-pass on the accel; they
sum to $1$, so it is exactly the complementary filter in disguise
(Figure~\ref{ch36:fig-compfreq}).

\begin{intuitionbox}
The cutoff between ``trust the gyro'' and ``trust the accel'' is set by $\alpha$ and the
sample time through the time constant $\tau = \dfrac{\alpha\,\Delta t}{1-\alpha}$.
Above $1/\tau$ the gyro wins; below it the accel wins. With $\Delta t = 0.01$\,s (a
$100$\,Hz loop) and $\alpha = 0.98$, $\tau \approx 0.49$\,s: trust the gyro for changes
faster than about half a second, trust the accel for anything slower. Want to lean
harder on the gyro (steadier, but drift creeps in)? Push $\alpha$ toward $1$.
\end{intuitionbox}

\begin{figure}[H]
\centering
\begin{tikzpicture}
\begin{axis}[width=11.5cm,height=5.0cm,
  xlabel={frequency, normalized to the crossover $f/f_c$},
  ylabel={magnitude}, xmin=0,xmax=5,ymin=0,ymax=1.15,
  domain=0:5, samples=160, legend pos=south east, legend cell align=left,
  grid=both, axis lines=left, every axis plot/.append style={thick}]
  % low-pass on accel: 1/sqrt(1+f^2)
  \addplot[targetgreen] {1/sqrt(1+x^2)};
  \addlegendentry{$|L|$: low-pass on accel}
  % high-pass on gyro: f/sqrt(1+f^2)
  \addplot[sourceblue] {x/sqrt(1+x^2)};
  \addlegendentry{$|H|$: high-pass on gyro}
  % sum of powers = 1: complementary
  \addplot[specred,dashed] {1};
  \addlegendentry{$|H|^2+|L|^2=1$}
\end{axis}
\end{tikzpicture}
\caption{The complementary pair in the frequency domain. Below the crossover $f_c$ the
low-pass (green) lets the accelerometer's slow, drift-free truth through; above it the
high-pass (blue) lets the gyro's fast detail through. The branches are designed so their
powers sum to one at every frequency (red dashed), so a steady true orientation passes
with unit gain and no notch. Sliding $\alpha$ slides $f_c$ left or right.}
\label{ch36:fig-compfreq}
\end{figure}

\begin{lstlisting}[language=Python,caption={Complementary filter for tilt from a 6-axis IMU.}]
import numpy as np

def complementary_tilt(gyro, accel, dt, alpha=0.98):
    """gyro: (N,3) rad/s, accel: (N,3) g. Returns roll,pitch (deg)."""
    roll = pitch = 0.0
    out = []
    for w, a in zip(gyro, accel):
        # accel-only absolute tilt (good when nearly still)
        roll_a  = np.degrees(np.arctan2(a[1], a[2]))
        pitch_a = np.degrees(np.arctan2(-a[0], np.hypot(a[1], a[2])))
        # gyro integration (good for fast motion), then blend
        roll  = alpha*(roll  + np.degrees(w[0])*dt) + (1-alpha)*roll_a
        pitch = alpha*(pitch + np.degrees(w[1])*dt) + (1-alpha)*pitch_a
        out.append((roll, pitch))
    return np.array(out)
\end{lstlisting}

\begin{figure}[H]
\centering
\begin{tikzpicture}
\begin{axis}[width=11.5cm,height=5.0cm,
  xlabel={time $t$ (s)}, ylabel={estimated pitch (deg)},
  xmin=0,xmax=10,ymin=-15,ymax=35, domain=0:10, samples=160,
  legend pos=north east, legend cell align=left, grid=both, axis lines=left]
  \addplot[targetgreen,thick,dashed] {20*(x>3)*(x<7)};
  \addlegendentry{true pitch}
  \addplot[specred,only marks,mark=*,mark size=0.5pt,samples=160]
     {20*(x>3)*(x<7) + 7*sin(deg(9*x)) + 5*cos(deg(13*x+1))};
  \addlegendentry{accel only (noisy)}
  \addplot[sourceblue,thick] {20*(x>3)*(x<7) + 6*exp(-2*x)};
  \addlegendentry{fused (complementary)}
\end{axis}
\end{tikzpicture}
\caption{A pitch step. Accel alone (red) is centered but jittery; integrated gyro alone
would track the step then drift. The complementary fusion (blue) follows the step
crisply and settles onto the true value with little noise.}
\label{ch36:fig-comp}
\end{figure}

\begin{applicationbox}
The complementary filter is everywhere because it is so cheap: self-balancing robots and
hoverboards, the early flight controllers of consumer drones (Multiwii, early
ArduPilot), camera gimbals, and countless hobby projects. On an 8-bit Arduino it is a
few multiplies per axis per loop. For two-axis tilt it is genuinely hard to beat for the
effort, and it is the right first thing to put on your wrist board.
\end{applicationbox}

\begin{connectionbox}
Look again at $\hat\theta = H(s)\theta_g + L(s)\theta_a$ with $H+L=1$. This is exactly
the complementary high-pass/low-pass pair from the filter-design chapters, and it is the
\emph{static} (constant-gain) ancestor of the Kalman filter, which lets the blend weights
\emph{vary in time} to be optimal given how noisy each sensor currently is.
\end{connectionbox}

\section{The Kalman filter as optimal fusion}
\label{ch36:kalman}

The complementary filter blends with a \emph{fixed} weight $\alpha$. But the right
weight is not fixed: when the device is whipping around, the accel is mostly motion and
should be trusted less; when it sits still, the accel is pure gravity and should be
trusted more. The Kalman filter is the principled answer to ``what is the
\emph{optimal}, time-varying blend, given everything I know about the noise?'' It is the
complementary filter with the weight recomputed every step from the statistics.

\begin{intuitionbox}
The Kalman filter keeps two things at all times: a best \emph{estimate} of the state
(here, orientation, maybe plus gyro bias) and a \emph{covariance} -- a number, or matrix,
saying how unsure it is. Each step has two beats. \textbf{Predict}: push the estimate
forward with the gyro, and \emph{grow} the uncertainty (we trusted a noisy gyro, so we
know less now). \textbf{Update}: when an absolute measurement arrives (accel/mag),
compare it to the prediction and step partway toward it -- and \emph{shrink} the
uncertainty. How far you step is the \emph{Kalman gain}, computed from the relative
sizes of your current uncertainty and the measurement's noise. It is a self-tuning
$\alpha$.
\end{intuitionbox}

Here is the whole idea in one dimension, which is enough to feel it. Let the state be a
scalar $x$ (an angle) with estimate $\hat x$ and variance $P$. The gyro gives a control
input $u=\omega\Delta t$ with process noise variance $Q$. The accel gives a measurement
$z$ of the angle with measurement noise variance $R$.

\begin{formalbox}
\textbf{Predict.}
\[
  \hat x^- = \hat x + u, \qquad P^- = P + Q .
\]
\textbf{Update} (Kalman gain $K$, then correct):
\[
  K = \frac{P^-}{P^- + R}, \qquad
  \hat x = \hat x^- + K\,(z - \hat x^-), \qquad
  P = (1-K)\,P^- .
\]
\end{formalbox}

Stare at the gain. $K = P^-/(P^-+R)$ lives between $0$ and $1$. If the measurement is
very noisy ($R$ large), $K\to 0$: ignore the accel, coast on the gyro. If our prediction
is very uncertain ($P^-$ large), $K\to 1$: snap to the accel. The correction
$\hat x = \hat x^- + K(z-\hat x^-)$ is \emph{identical in form} to the complementary
filter's $(1-\alpha)$ nudge toward the measurement -- but now the weight $K$ breathes
with the changing uncertainty.

\begin{whybox}
Why is this the \emph{optimal} blend, not just a reasonable one? If both noises are
Gaussian, ``predict'' and ``measure'' are two independent Gaussian opinions about the
same angle. The product of two Gaussians is another Gaussian whose mean is the
\emph{inverse-variance-weighted average} of the two means -- each opinion weighted by its
certainty. Working that average out gives exactly $K=P^-/(P^-+R)$. The Kalman filter is
nothing more than Bayes' rule for Gaussians, applied once per sample. That is the deep
link to the probability companion: the gain is the precision-weighting of two beliefs.
\end{whybox}

\begin{connectionbox}
Everything here is straight from the probability and linear-algebra books. The variance
$P$ is the spread of a belief; ``predict grows $P$, update shrinks $P$'' is uncertainty
propagation. The product-of-Gaussians giving a precision-weighted mean is Bayesian
updating with conjugate Gaussians. In full 3-D the scalars become vectors and matrices,
$z-\hat x^-$ becomes the \emph{innovation}, and $K$ becomes a matrix -- but the story is
unchanged: optimal fusion is inverse-variance weighting.
\end{connectionbox}

\begin{historybox}
Rudolf Kalman published the filter in 1960. Its first triumph was the Apollo program:
the guidance computer used a Kalman filter to fuse inertial measurements with star and
radar fixes, and it is no exaggeration that it helped get astronauts to the Moon and
back. Today it runs in GPS receivers, missiles, weather models, robot localization, and
the orientation chip in your phone. Stanley Schmidt at NASA Ames recognized its power and
pushed it into Apollo; the nonlinear variant is sometimes called the Schmidt--Kalman
filter.
\end{historybox}

\begin{trapbox}
Orientation is nonlinear (the quaternion kinematics, the trig in the accel model), so a
plain linear Kalman filter does not directly apply. Practitioners use the
\emph{extended} Kalman filter (EKF), which linearizes the model each step about the
current estimate, or the \emph{unscented} Kalman filter (UKF). These work well but are
fiddly: you must supply Jacobians (EKF), tune the noise matrices $Q$ and $R$, and guard
against divergence if the linearization is poor during violent motion. This fiddliness
is exactly why the lighter Madgwick and Mahony filters became popular.
\end{trapbox}

\section{The Madgwick and Mahony filters}
\label{ch36:madgwick-mahony}

A full quaternion EKF carries a $4\times4$ (or larger) covariance matrix and its matrix
algebra -- heavy for a tiny microcontroller. Around 2010 Sebastian Madgwick, then a PhD
student at Bristol, published a filter that achieves EKF-grade accuracy with a fraction
of the computation and \emph{one} intuitive tuning knob. It, and the closely related
Mahony filter, are what most open-source IMU libraries actually ship.

\begin{intuitionbox}
Both filters share one elegant idea, and it is a \emph{gradient-descent} idea you
already know from the calculus companion. Integrate the gyro to predict the quaternion.
Now ask: given my predicted orientation, where \emph{should} the accelerometer point if
the only force were gravity? Compare that predicted gravity direction to the measured
accel. The mismatch is an \emph{error}, and you can compute which way to nudge the
quaternion to shrink that error fastest -- the gradient. Take a small step downhill.
That small correction, blended with the gyro integration, is the whole filter: gyro for
speed, a gradient nudge toward gravity (and north, with a magnetometer) for the absolute
anchor.
\end{intuitionbox}

\begin{whybox}
Madgwick frames fusion as an optimization: find the quaternion $q$ that best aligns the
\emph{predicted} field directions (gravity, and magnetic north) with the
\emph{measured} accel and mag vectors. That is a least-squares problem,
\[
  \min_q \ \norm{\,f(q,\ \text{measured fields})\,}^2 ,
\]
and one step of gradient descent on it is cheap to write in closed form. The single knob
$\beta$ sets the descent step size -- the analog of the complementary $\alpha$ and the
Kalman $K$. Small $\beta$: trust the gyro more (smooth, slight drift). Large $\beta$:
pull harder toward accel/mag (drift-free, more noise). Mahony instead drives the same
mismatch through a PI controller and, as a bonus, its integral term \emph{estimates and
removes the gyro bias} -- a real practical win.
\end{whybox}

\begin{focusbox}
The unifying picture (Figure~\ref{ch36:fig-loop}): \emph{every} attitude filter in this
chapter is the same skeleton.
\begin{enumerate}
  \item \textbf{Predict} orientation by integrating the gyro (fast, drifts).
  \item Measure how far the prediction's gravity/north disagree with the accel/mag.
  \item \textbf{Correct} by nudging the prediction toward the absolute references.
\end{enumerate}
They differ only in \emph{how} they choose the correction strength:
\begin{itemize}
  \item Complementary: fixed weight $\alpha$.
  \item Kalman/EKF: optimal time-varying weight $K$ from the noise statistics.
  \item Madgwick: gradient-descent step of size $\beta$.
  \item Mahony: PI feedback with gain pair, plus online bias estimation.
\end{itemize}
\end{focusbox}

\begin{figure}[H]
\centering
\begin{tikzpicture}[>=Stealth,node distance=7mm,
  blk/.style={draw,rounded corners,align=center,minimum height=9mm,font=\small},
  every node/.style={font=\small}]
  \node[blk,fill=blue!5] (gyro) {gyro\\$\boldsymbol\omega$};
  \node[blk,fill=blue!5,right=18mm of gyro] (pred) {Predict\\integrate $\boldsymbol\omega$};
  \node[blk,fill=green!6,right=22mm of pred] (corr) {Correct\\nudge toward refs};
  \node[blk,right=18mm of corr] (out) {orientation\\$\hat q$};
  \node[blk,fill=red!5,below=12mm of corr] (acc) {accel/mag\\absolute refs};
  \draw[->] (gyro) -- (pred);
  \draw[->] (pred) -- node[above]{$\hat q^{-}$} (corr);
  \draw[->] (corr) -- (out);
  \draw[->] (acc) -- node[right]{error} (corr);
  % feedback loop
  \draw[->] (out.south) |- ($(out.south)+(0,-9mm)$) -| node[below,pos=0.25]{previous $\hat q$} (pred.south);
\end{tikzpicture}
\caption{Every attitude filter in this chapter is this loop. The gyro drives a
\emph{predict} step (fast, drifts); the accel/mag supply an absolute reference that the
\emph{correct} step nudges toward. The filters differ only in the correction strength:
fixed $\alpha$ (complementary), time-varying $K$ (Kalman), gradient step $\beta$
(Madgwick), or PI gains (Mahony).}
\label{ch36:fig-loop}
\end{figure}

\begin{table}[H]
\centering
\caption{Choosing an attitude filter.}
\label{ch36:tab-filters}
\begin{tabular}{@{}lllll@{}}
\toprule
Filter & State & Cost & Tuning & Best for \\
\midrule
Complementary & Euler/2 angles & tiny    & $\alpha$        & tilt only, 8-bit MCUs \\
Madgwick      & quaternion     & low     & $\beta$         & full 9-DoF on MCU \\
Mahony        & quaternion     & low     & $K_p,K_i$       & like Madgwick, removes bias \\
EKF           & quat ($+$ bias)& high    & $Q,R$ matrices  & best accuracy, more compute \\
\bottomrule
\end{tabular}
\end{table}

\begin{example}[One Kalman step by hand]
\label{ch36:ex-kalman}
A 1-D angle filter has $\hat x = 10.0^\circ$, $P = 4.0$ (deg$^2$). The gyro says the
angle rose by $u = 2.0^\circ$ this step with process noise $Q = 1.0$. The accel
measures $z = 13.5^\circ$ with noise $R = 9.0$.
\emph{Predict:} $\hat x^- = 12.0$, $P^- = 5.0$.
\emph{Gain:} $K = 5.0/(5.0+9.0) = 0.357$.
\emph{Update:} $\hat x = 12.0 + 0.357\,(13.5-12.0) = 12.54^\circ$, and
$P = (1-0.357)\cdot 5.0 = 3.21$. The estimate moved only a third of the way toward the
noisy accel (because $R>P^-$), and our uncertainty dropped from $5.0$ to $3.21$. Run it
again with the device held still and $P$ keeps shrinking; whip it around and $Q$ keeps
$P$ large so the filter leans on the accel. That breathing weight is the whole point.
\end{example}

\section{Getting a stable wrist orientation}
\label{ch36:wrist}

\begin{goalbox}
For the wrist controller, here is the recipe, in increasing order of effort.
\begin{enumerate}
  \item \textbf{Tilt mapping (start here).} Run a complementary filter on the 6-axis IMU
        for roll and pitch only. Map wrist tilt to two continuous MIDI controls -- e.g.
        pitch up/down to a pitch-bend, roll to a mod wheel. This is rock-solid, needs no
        magnetometer, and never drifts because gravity anchors it.
  \item \textbf{Full orientation.} Add the magnetometer and run Madgwick or Mahony to get
        a drift-resistant quaternion, hence heading too. Now ``which way the wrist
        points'' becomes a usable control, e.g. left/right pointing selects an
        instrument.
  \item \textbf{Bias and stillness.} Use Mahony's bias estimate (or a startup
        gyro-bias calibration while the arm rests) so the heading does not creep when the
        wrist is held still between notes.
\end{enumerate}
\end{goalbox}

\begin{buildbox}
A workable loop on an ESP32 or Teensy with an MPU-9250 / ICM-20948:
\begin{itemize}
  \item \textbf{Rate.} Run the filter at $100$--$200$\,Hz. Gyro integration error grows
        with $\Delta t$, so a steady, fast loop matters more than raw sensor quality.
  \item \textbf{Calibrate first.} Apply the gyro-bias, accel, and (especially)
        magnetometer hard/soft-iron calibration from the previous chapter \emph{before}
        fusion. The strongest magnetic distortion is your own board -- mount the IMU away
        from the battery, motors, and high-current traces.
  \item \textbf{Startup.} Hold the arm still for one second at boot; average the gyro to
        seed the bias and the accel to seed the initial quaternion. Tune $\beta$ (or
        $\alpha$) to taste: too low feels sluggish and drifts, too high feels twitchy.
  \item \textbf{Report.} Keep the quaternion as the working state; convert to Euler only
        at the moment you emit a MIDI value, and watch the $\pm 90^\circ$ pitch
        singularity when you do.
\end{itemize}
\end{buildbox}

\begin{lstlisting}[language=C,caption={Mahony-style correction core (per loop, quaternion q[4]).}]
// gx,gy,gz in rad/s; ax,ay,az normalized accel. Kp,Ki tuning gains.
// estimated direction of gravity from current quaternion:
float vx = 2*(q[1]*q[3] - q[0]*q[2]);
float vy = 2*(q[0]*q[1] + q[2]*q[3]);
float vz = q[0]*q[0] - q[1]*q[1] - q[2]*q[2] + q[3]*q[3];
// error = measured accel x estimated gravity (cross product):
float ex = (ay*vz - az*vy);
float ey = (az*vx - ax*vz);
float ez = (ax*vy - ay*vx);
ix += Ki*ex*dt; iy += Ki*ey*dt; iz += Ki*ez*dt;   // integral -> bias removal
gx += Kp*ex + ix;  gy += Kp*ey + iy;  gz += Kp*ez + iz;  // corrected rates
// integrate quaternion with corrected rates, then normalize q (omitted)
\end{lstlisting}

\begin{trapbox}
\textbf{What wrist orientation does, and where the fingers come from.} A wrist- or
forearm-worn IMU recovers the \emph{orientation of the forearm/wrist segment} -- tilt and
heading -- and discrete whole-hand gestures (a flick, a twist, a tap shock seen as an
accel spike). That alone gives the instrument several expressive axes. What the IMU
\emph{cannot} give you is which \emph{finger} moved: the IMU is bolted to the wrist, and
finger motion barely jostles it. But the fingers are far from out of reach. Each finger
is driven by tendons that run under the wrist and by forearm muscles whose electrical
activity reaches the skin, so finger movement is genuinely detectable from a band on the
wrist or forearm -- through \emph{force myography} (FMG, the tendon bulge under a band of
force sensors) or \emph{surface EMG} (the muscle's electrical pulses). With a short
per-finger calibration for each user, those signatures separate cleanly enough to drive
notes. That is the subject of the next chapter and the capstone. The one thing that
stays out of reach from the wrist is the \emph{continuous joint angle} of each finger --
how far each knuckle is bent at every instant. That information is not present in wrist
bulk motion and needs sensors \emph{on the fingers} (instrumented rings); see the
upgrade path in the capstone. Everything our instrument actually needs -- which finger,
how hard, when -- is on the doable side.
\end{trapbox}

\section*{Summary}
\label{ch36:summary}

Orientation is a single geometric object -- a rotation carrying the world frame onto the
body frame -- and choosing how to write it down matters. Euler angles are intuitive but
suffer gimbal lock; rotation matrices and quaternions are global and smooth, and the
unit quaternion (axis-and-half-angle in four numbers) is what real filters carry. No
single sensor gives orientation: the gyro is smooth but drifts (bias integrates to a
linear error), while the accelerometer and magnetometer are absolute but noisy and
motion-corrupted. Their flaws are complementary in frequency, so we fuse them. The
complementary filter is a high-pass on the gyro plus a low-pass on the accel with a fixed
weight; the Kalman filter makes that weight optimal and time-varying via inverse-variance
(Bayesian) weighting; Madgwick and Mahony reach the same goal with a cheap gradient or PI
correction toward gravity and north, Mahony also estimating gyro bias. For the wrist
controller, start with complementary tilt, graduate to Madgwick/Mahony for full
orientation -- and remember that the wrist IMU sees forearm pose and whole-hand gestures.
Which finger moved is recovered separately, from tendon (FMG) or muscle (EMG) sensing
with per-user calibration; only the continuous per-joint finger angle stays out of reach
of a wrist band.

\section{Exercises}
\label{ch36:exercises}

\begin{exercise}
Explain in your own words why three Euler angles cannot smoothly cover all orientations
without a singularity, and identify the orientation(s) at which the roll-pitch-yaw
convention gimbal-locks. What physical degree of freedom is lost there?
\end{exercise}

\begin{exercise}
A gyro has a constant bias of $0.5^\circ$/s and otherwise perfect output. Starting from a
known orientation, compute the heading error after $1$, $10$, and $60$ seconds of pure
integration. Now suppose the bias is zero but there is white rate noise; explain why the
error grows like $\sqrt{t}$ instead of $t$.
\end{exercise}

\begin{exercise}
Show that the discrete complementary filter
$\hat\theta_k = \alpha(\hat\theta_{k-1}+\omega_k\Delta t) + (1-\alpha)\theta_{a,k}$ is a
first-order high-pass on the gyro plus a first-order low-pass on the accel. Find the
crossover time constant $\tau$ and evaluate it for $\alpha=0.95$, $\Delta t=0.005$\,s.
\end{exercise}

\begin{exercise}
Why is the accelerometer trustworthy for orientation only when the device is nearly
still or averaged over time? Write the measured accel as gravity plus motion and explain
which part the low-pass branch of the complementary filter is trying to keep.
\end{exercise}

\begin{exercise}
Repeat Example~\ref{ch36:ex-kalman} but with the device held still, so you receive three
accel measurements $z = 13.5, 13.4, 13.6$ with no gyro motion ($u=0$). Run three
predict/update cycles and report how $\hat x$, $P$, and $K$ evolve. What happens to $K$
as $P$ shrinks, and what does that mean physically?
\end{exercise}

\begin{exercise}
Derive the Kalman gain $K=P^-/(P^-+R)$ from the product-of-two-Gaussians fact: if the
prediction is $\mathcal N(\hat x^-, P^-)$ and the measurement is $\mathcal N(z,R)$, show
the posterior mean is the inverse-variance-weighted average and rewrite it in the form
$\hat x^- + K(z-\hat x^-)$.
\end{exercise}

\begin{exercise}
Convert the unit quaternion $q=(\cos 30^\circ, 0, 0, \sin 30^\circ)$ to a rotation
matrix using the formula in Section~\ref{ch36:quat}. What axis and angle of rotation does
it represent? Verify the matrix is orthogonal with determinant $+1$.
\end{exercise}

\begin{exercise}
Both Madgwick and Mahony correct the gyro prediction toward gravity. In the C sketch, the
error is computed as the \emph{cross product} of measured accel and estimated gravity.
Explain geometrically why a cross product is the right error signal -- what its direction
and magnitude tell you about the misalignment, and why it vanishes when the two vectors
agree.
\end{exercise}

\begin{exercise}
Tune by experiment: implement the complementary filter and try $\alpha \in
\{0.90, 0.98, 0.999\}$ on the same recorded wrist motion. Describe qualitatively the
trade-off you see between responsiveness, noise, and drift, and state which you would
pick for (a) a fast percussive trigger and (b) a slow expressive filter sweep.
\end{exercise}

\begin{exercise}
Limits question. A friend proposes selling a single wristband that reconstructs every
finger \emph{joint angle} in real time from one 9-axis IMU. Explain precisely why a wrist
IMU cannot recover continuous joint angles, and what the IMU \emph{can} measure (forearm
pose, whole-hand gestures). Then explain which sensing modality you would add to the same
wrist band -- and why placing it there works -- to detect \emph{which finger moved}
without putting any hardware on the fingers themselves. Finally, state what you would have
to add on the fingers if you genuinely needed per-joint angles.
\end{exercise}
